%==============================================================================
% Appendix B: Computational checks and reproducibility
%==============================================================================
\section{Computational checks and reproducibility}
\label{app:B}

\subsection{Algebraic core of the $\chi$-bridge identity}
\label{app:B:bridge}

We record the algebraic core of the family-level identity
$\Ctop+9\chi=0$ of Section~\ref{sec:bridge:family}.  With the auxiliary
quantity $\Disc$ defined as
\begin{equation}
  \Disc = a^{2}-2ab-2ac+b^{2}-2bc+c^{2}+4u^{2}+4uv+4v^{2},
\end{equation}
the two quantities $\chi$ and $\Ctop$ are expressed as
\begin{equation}
  \chi = -\frac{\Disc}{12},
  \qquad
  \Ctop = \frac{3\Disc}{4}.
\end{equation}
Then
\begin{equation}
  \Ctop + 9\chi = \frac{3\Disc}{4} - \frac{9\Disc}{12} = 0.
\end{equation}
This identity holds for arbitrary values of $\Disc$ and does not assume any
specific topology.  The specialised values for the four topologies
(Table~\ref{tab:appA_topology}) are substitutions into this identity.

\subsection{Full topology branch taxonomy}
\label{app:B:branch_taxonomy}

The main admissibility classification across all 16 branches
($\Sthree/\Tthree/\Nilthree/\Solthree$ $\times$ $\EH/\AX/\VT/\MX$) is
given in Table~\ref{tab:appB1}.  Table~\ref{tab:four_topology_admissibility}
(Section~\ref{sec:classification}) is the summary of this classification
pattern.  The full $F$-branch expressions, the Hamiltonian constraints,
the auxiliary equations, and the auxiliary-shell solutions can be
reproduced with \code{admissibility_classification.py}.

\begin{table}[H]
\centering
\resizebox{\linewidth}{!}{%
\begin{tabular}{llcccl}
\toprule
Topology & Mode & $\chi$ & $F_{\mathrm{shell}}$ & Hessian determinant & Label \\
\midrule
$\Sthree$   & $\EH$ & $4$ & $4$ & n/a & \code{L_ADMISSIBLE} \\
$\Sthree$   & $\AX$ & $4$ & $4$ & $-2$ & \code{L_ADMISSIBLE} \\
$\Sthree$   & $\VT$ & $4$ & $4$ & $2q^{2}/9$ & \code{L_ADMISSIBLE} \\
$\Sthree$   & $\MX$ & $4$ & $4$ & $-4q^{2}(\kap^{4}\thetaNY^{2}+1)/9$ & \code{L_CONDITIONALLY_ADMISSIBLE} \\
$\Tthree$   & $\EH$ & $0$ & $0$ & n/a & \code{L_ADMISSIBLE} \\
$\Tthree$   & $\AX$ & $0$ & $0$ & $-2$ & \code{L_ADMISSIBLE} \\
$\Tthree$   & $\VT$ & $0$ & $0$ & $2q^{2}/9$ & \code{L_ADMISSIBLE} \\
$\Tthree$   & $\MX$ & $0$ & $0$ & $-4q^{2}(\kap^{4}\thetaNY^{2}+1)/9$ & \code{L_CONDITIONALLY_ADMISSIBLE} \\
$\Nilthree$ & $\EH$ & $-1/12$ & $-1/12$ & n/a & \code{L_ADMISSIBLE} \\
$\Nilthree$ & $\AX$ & $-1/12$ & $-1/12$ & $-2$ & \code{L_ADMISSIBLE} \\
$\Nilthree$ & $\VT$ & $-1/12$ & $-1/12$ & $2q^{2}/9$ & \code{L_ADMISSIBLE} \\
$\Nilthree$ & $\MX$ & $-1/12$ & $-1/12$ & $-4q^{2}(\kap^{4}\thetaNY^{2}+1)/9$ & \code{L_CONDITIONALLY_ADMISSIBLE} \\
$\Solthree$ & $\EH$ & $-1/3$ & $-1/3$ & n/a & \code{L_ADMISSIBLE} \\
$\Solthree$ & $\AX$ & $-1/3$ & $-1/3$ & $-2$ & \code{L_ADMISSIBLE} \\
$\Solthree$ & $\VT$ & $-1/3$ & $-1/3$ & $2q^{2}/9$ & \code{L_ADMISSIBLE} \\
$\Solthree$ & $\MX$ & $-1/3$ & $-1/3$ & $-4q^{2}(\kap^{4}\thetaNY^{2}+1)/9$ & \code{L_CONDITIONALLY_ADMISSIBLE} \\
\bottomrule
\end{tabular}%
}
\caption{Compact full branch taxonomy.}
\label{tab:appB1}
\end{table}

\subsection{Full orbit-sheet atlas}
\label{app:B:orbit_atlas}

The 24 branch/sheet rows of the reduced vacuum orbit atlas are given in
Table~\ref{tab:appB2}.  $\Sthree$ and $\Tthree$ contribute one row per
branch, while $\Nilthree$ and $\Solthree$ contribute two sheets
(expanding / collapsing).  Table~\ref{tab:orbit_class_counts}
(Section~\ref{sec:orbit:fullbranch}) summarises this branch/sheet count.
The orbit equations, constraint residuals, and numerical monitors of each
row can be reproduced with \code{orbit_atlas.py}.

\begin{table}[H]
\centering
\resizebox{\linewidth}{!}{%
\begin{tabular}{lllllc}
\toprule
Topology & Mode & Sheet & $\dot q_{0}$ & Orbit class & Check \\
\midrule
$\Sthree$   & $\EH$ & none       & no real solution & \code{NO_REAL_AUX_SHELL_ORBIT} & PASS \\
$\Sthree$   & $\AX$ & none       & no real solution & \code{NO_REAL_AUX_SHELL_ORBIT} & PASS \\
$\Sthree$   & $\VT$ & none       & no real solution & \code{NO_REAL_AUX_SHELL_ORBIT} & PASS \\
$\Sthree$   & $\MX$ & none       & no real solution & \code{NO_REAL_AUX_SHELL_ORBIT} & PASS \\
$\Tthree$   & $\EH$ & zero       & $0$              & \code{STATIC_OR_DEGENERATE}    & PASS \\
$\Tthree$   & $\AX$ & zero       & $0$              & \code{STATIC_OR_DEGENERATE}    & PASS \\
$\Tthree$   & $\VT$ & zero       & $0$              & \code{STATIC_OR_DEGENERATE}    & PASS \\
$\Tthree$   & $\MX$ & zero       & $0$              & \code{STATIC_OR_DEGENERATE}    & PASS \\
$\Nilthree$ & $\EH$ & expanding  & $\sqrt{3}/6$     & \code{MONOTONIC_EXPANSION}     & PASS \\
$\Nilthree$ & $\EH$ & collapsing & $-\sqrt{3}/6$    & \code{SINGULAR_APPROACH}       & PASS \\
$\Nilthree$ & $\AX$ & expanding  & $\sqrt{3}/6$     & \code{MONOTONIC_EXPANSION}     & PASS \\
$\Nilthree$ & $\AX$ & collapsing & $-\sqrt{3}/6$    & \code{SINGULAR_APPROACH}       & PASS \\
$\Nilthree$ & $\VT$ & expanding  & $\sqrt{3}/6$     & \code{MONOTONIC_EXPANSION}     & PASS \\
$\Nilthree$ & $\VT$ & collapsing & $-\sqrt{3}/6$    & \code{SINGULAR_APPROACH}       & PASS \\
$\Nilthree$ & $\MX$ & expanding  & $\sqrt{3}/6$     & \code{MONOTONIC_EXPANSION}     & PASS \\
$\Nilthree$ & $\MX$ & collapsing & $-\sqrt{3}/6$    & \code{SINGULAR_APPROACH}       & PASS \\
$\Solthree$ & $\EH$ & expanding  & $\sqrt{3}/3$     & \code{MONOTONIC_EXPANSION}     & PASS \\
$\Solthree$ & $\EH$ & collapsing & $-\sqrt{3}/3$    & \code{SINGULAR_APPROACH}       & PASS \\
$\Solthree$ & $\AX$ & expanding  & $\sqrt{3}/3$     & \code{MONOTONIC_EXPANSION}     & PASS \\
$\Solthree$ & $\AX$ & collapsing & $-\sqrt{3}/3$    & \code{SINGULAR_APPROACH}       & PASS \\
$\Solthree$ & $\VT$ & expanding  & $\sqrt{3}/3$     & \code{MONOTONIC_EXPANSION}     & PASS \\
$\Solthree$ & $\VT$ & collapsing & $-\sqrt{3}/3$    & \code{SINGULAR_APPROACH}       & PASS \\
$\Solthree$ & $\MX$ & expanding  & $\sqrt{3}/3$     & \code{MONOTONIC_EXPANSION}     & PASS \\
$\Solthree$ & $\MX$ & collapsing & $-\sqrt{3}/3$    & \code{SINGULAR_APPROACH}       & PASS \\
\bottomrule
\end{tabular}%
}
\caption{Compact full orbit-sheet atlas.}
\label{tab:appB2}
\end{table}

\subsection{Summary of main computational checks}
\label{app:B:checks}

Table~\ref{tab:appB3} lists the logical verification targets.  The first
two checks are evaluated by the same symbolic script, because the
family-level identity and its topology specialisations share the same
derived $(a,b,c,u,v)$ data.

\begin{table}[H]
\centering
\begingroup
\renewcommand{\arraystretch}{2.0}
\begin{tabularx}{\linewidth}{YYYY}
\toprule
Check & Content & \shortstack{Related\\section} & Script \\
\midrule
$\chi$-bridge derivation
  & Extract $(a,b,c,u,v)$ for $\Sthree/\Tthree/\Nilthree/\Solthree$,
    derive $\chi$ from the LC spatial curvature, compute the raw $\PintMX$
    from the EC+NY MX curvature, and extract $\Ctop$
  & Section~\ref{sec:pchannel:Pint}, \ref{sec:bridge:family}
  & \code{chi_bridge_symbolic.py} \\
$\Pform$ cancellation
  & Construct the curvature with the EC+NY engine for each torsion branch
    $\AX/\VT/\MX$ and compute $\Pform$
  & Section~\ref{sec:pchannel:Pform}
  & \code{pform_cancellation.py} \\
admissibility classification
  & Derive the lapse constraint and the auxiliary closure symbolically from
    the engine-derived $F_{\mathrm{branch}}$
  & Section~\ref{sec:hamiltonian}, \ref{sec:classification}
  & \code{admissibility_classification.py} \\
orbit atlas
  & Derive the orbit equation and the orbit class from the auxiliary
    shell of the engine-derived $F_{\mathrm{branch}}$
  & Section~\ref{sec:orbit}
  & \code{orbit_atlas.py} \\
\bottomrule
\end{tabularx}
\endgroup
\caption{Computational checks and their roles.}
\label{tab:appB3}
\end{table}
